\documentclass[10pt,a4paper]{amsart}
\usepackage[margin=19mm]{geometry}
\usepackage{amsmath,amssymb,mathtools}
\usepackage[scr=rsfs,cal=euler]{mathalfa}
\usepackage{xfrac}
\usepackage[unicode,hidelinks,bookmarks=false]{hyperref}
\newcommand{\ie}{i.e.\ }
\newcommand{\cf}{cf.\ }
\newcommand{\resp}{respectively\ }
\newcommand{\Subsection}{\subsection}
\providecommand{\nobreakdash}{\nobreak\hbox{-}}
\setlength{\emergencystretch}{2em}
\multlinegap0pt
\usepackage{xcolor}
\usepackage{amssymb}
\usepackage{enumerate}
\usepackage{bm}
\usepackage{tikz}
\newtheorem{theorem}{Theorem}
\newtheorem{lemma}{Lemma}
\newcommand{\NN}{\mathbb{N}}
\newcommand{\Pra}[1]{\mathbb{P}\left(#1\right)}
\newcommand{\bB}{\mathbf{B}}
\newcommand{\bN}{\mathbb{N}}
\newcommand{\bS}{\mathbf{S}}
\newcommand{\cA}{\mathcal{A}}
\newcommand{\cB}{\mathcal{B}}
\newcommand{\cS}{\mathcal{S}}
\renewcommand{\ge}{\geqslant}
\renewcommand{\le}{\leqslant}
\newcommand{\product}{\Lambda}
\newcommand{\productO}{\Lambda(\cS)}
\title{Arrangements of Consecutive Numbers in Mallows Permutations}
\author{Katarzyna Rybarczyk}
\date{Source: arXiv:2606.12410v1 (CC BY 4.0)}
\begin{document}
\maketitle

\noindent\emph{Source and licence.} Arrangements of Consecutive Numbers in Mallows Permutations. Version arXiv:2606.12410v1 (CC BY 4.0), distributed by arXiv under the Creative Commons Attribution 4.0 International licence, http://creativecommons.org/licenses/by/4.0/, as recorded in the arXiv licence metadata for this exact version and shown on the abstract page https://arxiv.org/abs/2606.12410v1. Full licence notice: \texttt{licenses/arxiv-2606.12410-CC-BY-4.0.txt}.\par\smallskip

\noindent\textit{Editorial scope.} Counted unit: the article's theorem Thm:Qk:value (source0.tex lines 374--393). The article's complete original proof of it is delivered with it (source0.tex lines 848--933, one \texttt{proof} environment of 86 lines, which the article places away from the statement). No mathematics is written, repaired or re-derived: the statement and the proof are the article's own lines, in the article's own notation, and no retained line is substituted or re-flowed.  Retained uncounted as the source-local context the proof needs: lines 241--243; lines 731--734.  Nothing was omitted from any retained span, so the package carries no omission record.  No retained line contains a citation, so the wrapper carries no bibliography and no bibliography entry.  No retained line contains a figure environment, an image-inclusion command or drawing code, so no image is delivered and the package declares no asset dependency.  Editorial formatting only, in addition: an AMS \texttt{amsart} preamble that restates the article's own declarations and macros used by the retained lines, namely the environments \texttt{theorem}, \texttt{lemma} on separate counters, together with the standard packages those lines need.  The retained environments print in this document as Theorem 1, Lemma 1 in order of appearance, whereas the article prints the same items with the numbering of its own full document.  The complete native source of the article as distributed in the arXiv e-print for this exact version is bundled unchanged as \texttt{sources/202622/source0.tex}.
\begin{equation}\label{Eq:Def:Product}
	\product(\bS_\ell)=\prod_{j=1}^{\ell}(1-q^j).
\end{equation}

\begin{lemma}\label{Lem:AuxiliarySum} Let $q =q(n)\in (0,1)$ and $k\ge \ell -\log_q n$. Then, 
	\[\sum_{j=0}^{k-\ell+1}\left(q^{\ell}\right)^j\prod_{r=1}^{\ell-1}(1-q^{j+r})
	=(1+O(n^{-\ell}))\frac{1}{1-q^\ell}\prod_{i=1}^{\ell-1}\frac{(1-q^i)}{(1-q^{\ell+i})}.\]
\end{lemma}

\begin{theorem}\label{Thm:Qk:value}
	Let $\ell\ge 2$ be a constant, $\emptyset\neq\cS\subseteq \bS_\ell$, $q=q(n)\in (0,1)$, and $t=t(n)\in \NN$ be such that $(1-q)t/\ln n\to \infty$. Then for $k=k(n)\ge \ell -\log_q n$
	\begin{equation*}
		\begin{split}
			\bB_{k}
			=(1+O(n^{-1}))
			\frac{\productO}{1-q^\ell}
			\prod_{i=1}^{\ell-1}\frac{(1-q^i)}{(1-q^{\ell+i})}+O(\productO n^{-\ell}),
		\end{split}
	\end{equation*}
	where $\productO$ is defined in \eqref{Eq:Def:Product}.
	
	Moreover for $1-q=o(1)$ we have
	\begin{equation*}
		\begin{split}
			\bB_{k}=(1+O(n^{-1}+O(1-q)))\frac{\productO}{1-q}\frac{(\ell-1)!}{(2\ell-1)!}
			=(1+o(1))\frac{|\cS|[(\ell-1)!]^2}{(2\ell-1)!}(1-q)^{\ell-1}.
		\end{split}
	\end{equation*}
\end{theorem}

\begin{proof}[Proof of Theorem~\ref{Thm:Qk:value}]
	
	Events $\cA[0,i]$, $i=0,1,2,\ldots$, are mutually exclusive and $\cB[0]=\bigcup_{i\in \NN_0}\cA[0,i]$. Therefore
	\begin{multline}\label{Eq:Q0}
		\bB_{0}=\Pra{\cB[0]}=\sum_{i=0}^{\infty}\Pra{\cA[0,i]}=\\=\sum_{i=0}^{\infty}\left(\prod_{j=1}^{i}\Pra{Z_j>\ell}\right)\Pra{\cA[0,0]}
		=\sum_{i=0}^{\infty}(q^\ell)^i\Pra{\cA[0,0]}=\frac{\productO}{1-q^\ell},
	\end{multline}
	where $i$ represents the number of times that $Z_j>\ell$ before in the Mallows infinite process one chooses the first number in $[\ell]$. By convention we set $\prod_{j=1}^0a_j=1$.
	
	Now let us assume that $k\ge 1$.
	Before we study a general formula for $\bB_{k}$ we start with the value of $\Pra{\cA[k,0]}$. By the properties of the geometric distribution $\operatorname{Geom}(1-q)$ we have 
	\begin{equation}\label{Eq:Qk:Ak0}
		\begin{split}
			\Pra{\cA[k,0]}
			&=\sum_{\sigma\in \cS}\prod_{j=1}^{\ell}\Pra{Z_j=k+\sigma(j)-|\sigma([j-1])\cap [\sigma(j)-1]|}\\
			&=\sum_{\sigma\in \cS}\prod_{j=1}^{\ell}q^k\Pra{Z_j=\sigma(j)-|\sigma([j-1])\cap [\sigma(j)-1]|}\\
			&=q^{k\ell}\productO.
		\end{split}
	\end{equation}
	
	Now, we determine the recursive formula for $\bB_{k}$. As before, $i$ counts the number of times the geometric random variable $Z_j$, $j=1,2,\ldots$, in the infinite Mallows process, takes values greater than $k+\ell$, before it takes a value at most $k+\ell$.
	\begin{align}\nonumber
		\bB_{k}
		&=\sum_{i=0}^{\infty}\left(\prod_{j=1}^{i}\Pra{Z_j>k+\ell}\right)\left(\Pra{\cA[k,0]}+\Pra{Z_{i+1}\le k}\bB_{k-1}\right)\\
		\nonumber
		&=\sum_{i=0}^{\infty}(q^{\ell+k})^i\left(\Pra{\cA[k,0]}+(1-q^k)\bB_{k-1,l}\right)\\
		\label{Eq:Qk:recurrence}
		&=\frac{q^{k\ell}\productO}{1-q^{\ell+k}}+\frac{(1-q^k)}{1-q^{k+\ell}}\bB_{k-1},
	\end{align} 
	where we used \eqref{Eq:Qk:Ak0} in the last line.
	By applying \eqref{Eq:Qk:recurrence} one can easily show by induction that, for any $k\in \bN_0$,
	\begin{equation}\label{Eq:Qk:Sum0}
		\bB_{k}=\sum_{i=0}^k q^{i\ell}\frac{\product[i+1,k]\productO}{\product[\ell+i,\ell+k]}.
	\end{equation}
	
	Now we split the sum \eqref{Eq:Qk:Sum0} into two. 
	\begin{equation}\label{Eq:Qk:Sum01}
		\bB_{k}=\sum_{i=0}^{k-\ell+1} q^{i\ell}\frac{\product[i+1,k]\productO}{\product[\ell+i,\ell+k]}+\sum_{i=k-\ell+2}^{k} q^{i\ell}\frac{\product[i+1,k]\productO}{\product[\ell+i,\ell+k]}.
	\end{equation}
	
	
	The first sum in \eqref{Eq:Qk:Sum01} contains the terms with $i\le k-\ell+1$. For $i\le k-\ell+1$ we have
	\begin{equation}\label{Eq:Qk:cancelation}
		\frac{\product[i+1,k]}{\product[\ell+i,\ell+k]}
		=
		\frac{\product[i+1,i+\ell-1]\product[i+\ell,k]}{\product[i+\ell,k]\product[k+1,\ell+k]}
		=\frac{\product[i+1,i+\ell-1]}{\product[k+1,\ell+k]}.
	\end{equation}
	Moreover, we consider $k\ge \ell -\log_q n$. For them
	\[
	q^k\le q^{\ell-\log_q n}\le n^{-1}.\]
	Therefore
	\begin{equation}\label{Eq:Qk:bottomtems1}  
		1\ge\product[k+1,\ell+k] = \prod_{j=1}^{\ell}(1-q^{k+i}) \ge (1-q^{k})^\ell \ge (1-n^{-1})^\ell\ge 1-\ell n^{-1}=1+O(n^{-1}).
	\end{equation}
	Similarly, for the second sum in \eqref{Eq:Qk:Sum01} we have $i=k-\ell+2,\ldots,k$.  Using \eqref{Eq:Qk:bottomtems1} 
	\begin{equation}\label{Eq:Qk:bottomtems2}  
		1\ge \product[\ell+i,\ell+k]\ge \product[k+1,\ell+k] = 1+O(n^{-1}).
	\end{equation}
	In addition, for $i=k-\ell+2,\ldots,k$, since $k\ge \ell -\log_q n$,
	\begin{equation}\label{Eq:Qk:topterms}
		q^{i\ell}\le q^{(k-\ell)\ell}\le n^{-\ell}=O(n^{-\ell}). 
	\end{equation}
	And obviously $\product[\ell+i,\ell+k]\le 1$. Therefore
	we apply \eqref{Eq:Qk:cancelation} and \eqref{Eq:Qk:bottomtems1} for the first sum and \eqref{Eq:Qk:bottomtems2} and \eqref{Eq:Qk:topterms} for the second sum in \eqref{Eq:Qk:Sum01} to obtain
	\begin{equation}\label{Eq:Qk:Sum1}
		\begin{split}
			\bB_{k}
			&=\productO\sum_{i=0}^{k-\ell+1}\left(q^{\ell}\right)^i\frac{\product[i+1,i+\ell-1]}{\product[k+1,k+\ell]}+\sum_{i=k-\ell+2}^{k}q^{i\ell}\frac{\product[i+1,k]\productO}{\product[\ell+i,\ell+k]}\\
			&=(1+O(n^{-1}))\productO\sum_{i=0}^{k-\ell+1}\left(q^{\ell}\right)^i\product[i+1,i+\ell-1]+O(n^{-\ell}\productO)\\
			&=(1+O(n^{-1}))\productO\left(\sum_{j=0}^{k-\ell+1}\left(q^{\ell}\right)^j\prod_{r=1}^{\ell-1}(1-q^{j+r})\right)+O(n^{-\ell}\productO).
		\end{split}
	\end{equation}
	The first part of Theorem~\ref{Thm:Qk:value} follows immediately from Lemma~\ref{Lem:AuxiliarySum} applied to the above equation.
	
	Now we assume that $1-q=o(1)$. Then for any $j=O(1)$ we have
	\[
	1-q^j=1-[1-(1-q)]^j=j(1-q)(1+O(1-q)).
	\]
	Therefore
	\begin{align}\nonumber
		\bB_{k}
		&=(1+O(n^{-1}))\frac{\productO}{1-q^\ell}\prod_{i=1}^{\ell-1}\frac{(1-q^i)}{(1-q^{\ell+i})}+O(\productO n^{-\ell})\\	\nonumber		&=(1+O(n^{-1}+O(1-q)))\frac{\productO}{\ell(1-q)}\frac{(\ell-1)!\ell!}{(2\ell-1)!}+O(\productO n^{-\ell})\\
		\label{Eq:q:tendsto1}			&=(1+O(n^{-1}+O(1-q)))\frac{\productO}{1-q}\frac{[(\ell-1)!]^2}{(2\ell-1)!}.
	\end{align}
\end{proof}

\end{document}
